\section{Datasets and Experimental Setup}
\label{sec:setup}

Every study in this report trains a model on a short observed window of one
trajectory and then scores it on a later window it never saw. This section
fixes, once, what those trajectories are, how they are generated and split,
and which model and optimizer settings each system is trained with, so that
later sections only need to state what they change.

\subsection{The four dynamical systems}

Three of the four systems are defined by a known ordinary differential
equation; the fourth is deliberately \emph{not}, so that it can test whether
fixes derived from a known equation still work when no such equation exists.

\paragraph{Lotka-Volterra predator-prey.} The base paper's benchmark
\cite{koenig2024kanode}: prey $u_1$ and predators $u_2$ follow
\begin{equation}
\dot u_1 = \alpha u_1 - \beta u_1 u_2, \qquad
\dot u_2 = \delta u_1 u_2 - \gamma u_2,
\qquad (\alpha,\beta,\gamma,\delta)=(1.5,\,1.0,\,3.0,\,1.0),
\label{eq:lv}
\end{equation}
from $\mathbf{u}_0=(1,1)$. The system has a conserved quantity,
$V(u_1,u_2)=\delta u_1-\gamma\ln u_1+\beta u_2-\alpha\ln u_2$, so every
trajectory is a closed orbit. From this initial state the period is $3.32$
time units, so the training window $[0,3.5]$ shows the model almost exactly
one cycle. An accurate model must therefore keep returning to the
same orbit, and any error in the learned field shows up as a slow drift in
amplitude or phase that grows with the extrapolation horizon. (The base paper
writes the same system with the names of $\gamma$ and $\delta$ swapped; the
dynamics are identical.)

\paragraph{Damped pendulum.} Angle $\theta$ and angular velocity $\omega$,
\begin{equation}
\dot\theta=\omega, \qquad \dot\omega=-\tfrac{g}{L}\sin\theta-\mu\,\omega,
\qquad g/L=9.81,\ \mu=0.5,
\label{eq:pendulum}
\end{equation}
from $(\theta_0,\omega_0)=(2.0,\,0)$ rad, a large initial swing where
$\sin\theta\approx\theta$ is a poor approximation, so the model has to learn
a genuinely nonlinear restoring force. Friction makes the trajectory spiral
into the equilibrium $(0,0)$: the mechanical energy
$E=\tfrac12 mL^2\omega^2+mgL(1-\cos\theta)$ satisfies
$\dd E/\dd t=-mL^2\mu\,\omega^2\le0$, so it can never increase. A model can
fit the trajectory closely and still break this law, which makes it an
independent physical check (Section~\ref{sec:phase2}). The damping study (Section~\ref{sec:track-d})
varies $\mu\in\{0.1,0.5,1,2,5,8\}$.

\paragraph{SIR epidemic.} Susceptible, infected and recovered fractions of a
fixed population,
\begin{equation}
\dot S=-\beta SI, \qquad \dot I=\beta SI-\gamma I, \qquad \dot R=\gamma I,
\qquad \beta=0.35,\ \gamma=0.10\ \ (R_0=\beta/\gamma=3.5),
\label{eq:sir}
\end{equation}
from $(S_0,I_0,R_0)=(0.99,\,0.01,\,0)$, with time measured in days. The
right-hand sides sum to zero, so $S+I+R=1$ holds exactly for all time. The
infected curve rises to a single peak and then decays, so the extrapolation
window tests whether the model has learned the \emph{end} of an outbreak from
data that mostly shows its growth.

\paragraph{Synthetic non-mechanistic epidemic curve.} To test whether the SIR
fixes transfer to data that no compartmental ODE generated
(Section~\ref{sec:track-e}), an infection curve is built directly from a
closed-form skewed bell shape rather than from any differential equation,
\begin{equation}
\tilde I(t)=\exp\!\Big(-\tfrac12\big(\tfrac{t-38}{14}\big)^2\Big)\Big(1+\tanh\big(1.4\,\tfrac{t-38}{14}\big)\Big)+\varepsilon_t,
\qquad \varepsilon_t\sim\mathcal N(0,\,0.01^2),
\label{eq:synthetic-epidemic}
\end{equation}
sampled daily for $t\in\{0,\dots,120\}$, clipped at zero and smoothed with a
7-day moving average, the way public case counts are usually reported. The
second state is a cumulative-recovered proxy,
$\tilde R(t)=0.08\sum_{s\le t}\tilde I(s)$. Both columns are min-max
normalized to $[0,1]$. The $\tanh$ factor makes the rise steeper than the
decay, which no symmetric SIR peak reproduces exactly. This curve is
\emph{synthetic}: it imitates the shape of a real outbreak but is not
recorded case data.

\subsection{Ground truth, sampling and train/test splits}

The three mechanistic systems are integrated to near machine precision
(SciPy's \texttt{solve\_ivp}, eighth-order Dormand-Prince, relative and
absolute tolerance $10^{-12}$), far tighter than any error the learned models
reach, so the reference trajectory contributes no measurable error of its
own. Each trajectory is sampled on a uniform grid of spacing $\Delta t$ and
cut at $t_{\text{train}}$: points up to and including $t_{\text{train}}$ form
the training set, and every later point forms the test (extrapolation) set.
There is one trajectory per system, trained full-batch; there is no random
shuffling and no validation split, because the test window is defined by
time, not sampled from the same distribution as training.

\begin{table}[!ht]
\centering
\footnotesize
\renewcommand{\arraystretch}{1.15}
\caption{Dataset definitions. ``Train'' is $[0,t_{\text{train}}]$; ``Test''
is $(t_{\text{train}},t_{\text{end}}]$. Point counts include both endpoints
of the training window.}
\label{tab:datasets}
\begin{tabularx}{\textwidth}{@{}L{2.6cm} C{1.25cm} C{1.3cm} C{1.3cm} C{1.9cm} C{1.9cm} X@{}}
\toprule
\hdrrow \thh{System} & \thh{State dim.} & \thh{Spacing $\Delta t$} & \thh{$t_{\text{end}}$} & \thh{Train window (points)} & \thh{Test window (points)} & \thh{Test window covers}\\
\midrule
Lotka-Volterra & 2 & 0.1 & 14 & $[0,3.5]$ (36) & $(3.5,14]$ (105) & $\approx3.2$ unseen cycles; $(14,28]$ adds $\approx4.2$ more (Table~\ref{tab:extrap-horizon})\\
Damped pendulum & 2 & 0.05 & 10 & $[0,5]$ (101)\textsuperscript{a} & $(5,10]$ (100) & late decay toward rest\\
SIR & 3 & 0.5 & 80 & $[0,50]$ (101) & $(50,80]$ (60) & end of the outbreak\\
Synthetic epidemic & 2 & 1 day & 120 & days 0--44 (45)\textsuperscript{b} & days 45--120 (76) & decay after the peak\\
\bottomrule
\end{tabularx}

\smallskip
\raggedright\footnotesize
\textsuperscript{a}The original recipe used $[0,3]$ (61 points); Section~\ref{sec:phase2} explains why it was extended.
\textsuperscript{b}Split controls in Section~\ref{sec:track-e} move the cut to day 60 and day 70.
\end{table}

\paragraph{Observation noise.} Where a study adds noise (the Lotka-Volterra
noise sweep, Section~\ref{sec:track-e}), independent Gaussian noise
$\mathcal N(0,\sigma^2)$ with $\sigma\in\{0,0.01,0.05,0.1\}$ is added to the
\emph{training} observations only. The test window is always scored against
the clean reference trajectory, so the reported error measures how well the
model recovered the true dynamics, not how well it reproduced the noise.

\subsection{Models and training settings}

All KAN-ODEs use the edge parameterization of
Equation~\ref{eq:kan-edge} with a $\tanh$ input normalizer and a SiLU
residual path (Section~\ref{sec:numerics} gives the exact computation). A
KAN layer from $n_{\text{in}}$ to $n_{\text{out}}$ nodes with $G$ basis
functions per edge has
\begin{equation}
P_{\text{layer}} = n_{\text{out}}\,n_{\text{in}}\,(G+1)
\label{eq:param-count}
\end{equation}
parameters: $G$ basis coefficients plus one residual weight on each of its
$n_{\text{in}}n_{\text{out}}$ edges. The size of the network differs by
system, because SIR's three-dimensional state and the pendulum's
sharper nonlinearity needed more capacity than Lotka-Volterra.

\begin{table}[!ht]
\centering
\footnotesize
\renewcommand{\arraystretch}{1.15}
\caption{Model and training settings for each system's main runs. Studies
that deviate (solver or basis ablations, the damping sweep's 2{,}000-epoch
budget and its range of damping coefficients) say so where they are
reported.}
\label{tab:train-settings}
\begin{tabularx}{\textwidth}{@{}L{3.3cm} C{2.5cm} C{2.3cm} C{2.3cm} X@{}}
\toprule
\hdrrow \thh{Setting} & \thh{Lotka-Volterra} & \thh{Pendulum} & \thh{SIR} & \thh{Synthetic epidemic}\\
\midrule
KAN widths & $[2,10,2]$ & $[2,10,2]$ & $[3,16,3]$ & $[2,10,2]$\\
Basis, grid size $G$ & RBF, 5 & RBF, 8 & RBF, 8 & RBF, 5\\
Parameters (Eq.~\ref{eq:param-count}) & 240 & 360 & 864 & 240\\
Solver, substeps $m$ & Tsit5, 2 & Tsit5, 2 & Tsit5, 2 & Tsit5, 2\\
Learning rate (Adam) & $2\times10^{-3}$ & $3\times10^{-3}$ & $3\times10^{-3}$ & $2\times10^{-3}$\\
Gradient-norm clipping & none (1.0 in the extended-budget reruns, Section~\ref{sec:phase4}) & 1.0 & 1.0 & 1.0\\
Epochs & 10{,}000 & 10{,}000 & 10{,}000 & 2{,}000 (probes), 5{,}000 (full)\\
Structural fixes & none & none & time scale $s=10$, projection, gate on $I$ & time scale $s=24$; projection and gate tested as arms\\
Seed & 42 & 42 & 42 & 42\\
\bottomrule
\end{tabularx}
\end{table}

\paragraph{MLP baselines.} Two multilayer-perceptron Neural ODEs are trained
on Lotka-Volterra under exactly the same protocol for comparison: the base
paper's literal architecture, $[2,50,2]$ with $\tanh$ (252 parameters), and a
parameter-matched deeper network, $[2,14,8,8,2]$ with SiLU (252 parameters),
introduced because the literal one fails to train under this protocol
(Table~\ref{tab:mlp-breakdown}).

\paragraph{Why 10{,}000 epochs.} The base paper trains for $10^5$ epochs.
This project fixed its budget at $10^4$ for three reasons. First, cost: the
full ablation is 24 Lotka-Volterra configurations plus three other systems,
and at the measured $0.03$--$0.67$~s per epoch on a CPU a $10^5$-epoch
sweep of the Lotka-Volterra ablation alone would take on the order of a week
of serial CPU time, before any other system or follow-up study. Second, the paper itself reports that its
KAN-ODE reaches $2.6\times10^{-5}$ training loss within $10^4$ epochs,
already below its MLP's $10^5$-epoch result, so $10^4$ is where the paper's
own speed claim is made. Third, the ablations compare configurations against
each other at a common budget, so what matters is whether the
\emph{ranking} is stable, not whether each run has fully converged. That last
assumption was tested directly by retraining five configurations to
25{,}000--50{,}000 epochs; Section~\ref{sec:phase4} reports the result, which
confirms the ranking within the KAN variants but reverses the KAN-versus-MLP
comparison.

\paragraph{Training objective and model selection.} Each epoch integrates the
learned field once over the whole training window from the known initial
state and minimizes the mean squared error against the observed points,
\begin{equation}
\mathcal L(\theta)=\frac{1}{N_{\text{train}}\,d}\sum_{n=0}^{N_{\text{train}}-1}\big\lVert \hat{\mathbf u}_\theta(t_n)-\mathbf u(t_n)\big\rVert_2^2,
\label{eq:train-loss}
\end{equation}
where $d$ is the state dimension and $\hat{\mathbf u}_\theta(t_n)$ is the
integrator's prediction. No sparsity or smoothness regularization is added
in any run. Adam uses a constant learning rate with no schedule, matching the
base paper. The checkpoint that is reported is the one with the lowest
\emph{training} loss seen during the run; the test window is monitored but
never used to select a model, so no reported number is tuned on the data it
is scored on. Algorithm~\ref{alg:train} in Section~\ref{sec:numerics} gives
the complete loop.

\FloatBarrier
